% Options for packages loaded elsewhere
\PassOptionsToPackage{unicode}{hyperref}
\PassOptionsToPackage{hyphens}{url}
\PassOptionsToPackage{dvipsnames,svgnames,x11names}{xcolor}
%
\documentclass[
  11pt,
]{article}
\usepackage{amsmath,amssymb}
\usepackage{lmodern}
\usepackage{iftex}
\ifPDFTeX
  \usepackage[T1]{fontenc}
  \usepackage[utf8]{inputenc}
  \usepackage{textcomp} % provide euro and other symbols
\else % if luatex or xetex
  \usepackage{unicode-math}
  \defaultfontfeatures{Scale=MatchLowercase}
  \defaultfontfeatures[\rmfamily]{Ligatures=TeX,Scale=1}
\fi
% Use upquote if available, for straight quotes in verbatim environments
\IfFileExists{upquote.sty}{\usepackage{upquote}}{}
\IfFileExists{microtype.sty}{% use microtype if available
  \usepackage[]{microtype}
  \UseMicrotypeSet[protrusion]{basicmath} % disable protrusion for tt fonts
}{}
\makeatletter
\@ifundefined{KOMAClassName}{% if non-KOMA class
  \IfFileExists{parskip.sty}{%
    \usepackage{parskip}
  }{% else
    \setlength{\parindent}{0pt}
    \setlength{\parskip}{6pt plus 2pt minus 1pt}}
}{% if KOMA class
  \KOMAoptions{parskip=half}}
\makeatother
\usepackage{xcolor}
\IfFileExists{xurl.sty}{\usepackage{xurl}}{} % add URL line breaks if available
\IfFileExists{bookmark.sty}{\usepackage{bookmark}}{\usepackage{hyperref}}
\hypersetup{
  pdftitle={Project Proposal (Week 2), Ventilation Fault Detection, Diagnosis and Graceful Degradation},
  pdfauthor={Wassim Akachi, Project group 5},
  colorlinks=true,
  linkcolor={RoyalBlue},
  filecolor={Maroon},
  citecolor={Blue},
  urlcolor={Blue},
  pdfcreator={LaTeX}}
\urlstyle{same} % disable monospaced font for URLs
\usepackage[a4paper,margin=2.2cm,top=1.9cm,bottom=1.7cm]{geometry}
\usepackage{longtable,booktabs,array}
\usepackage{calc} % for calculating minipage widths
% Correct order of tables after \paragraph or \subparagraph
\usepackage{etoolbox}
\makeatletter
\patchcmd\longtable{\par}{\if@noskipsec\mbox{}\fi\par}{}{}
\makeatother
% Allow footnotes in longtable head/foot
\IfFileExists{footnotehyper.sty}{\usepackage{footnotehyper}}{\usepackage{footnote}}
\makesavenoteenv{longtable}
\setlength{\emergencystretch}{3em} % prevent overfull lines
\providecommand{\tightlist}{%
  \setlength{\itemsep}{0pt}\setlength{\parskip}{0pt}}
\setcounter{secnumdepth}{-\maxdimen} % remove section numbering
\usepackage{newunicodechar}
\usepackage{textcomp}
\usepackage{amssymb}
\newunicodechar{°}{\textdegree}
\newunicodechar{₂}{\textsubscript{2}}
\newunicodechar{→}{\ensuremath{\rightarrow}}
\newunicodechar{≤}{\ensuremath{\leq}}
\newunicodechar{≥}{\ensuremath{\geq}}
\newunicodechar{≈}{\ensuremath{\approx}}
\newunicodechar{✓}{\ensuremath{\checkmark}}
\newunicodechar{–}{--}
\newunicodechar{—}{---}
\usepackage{etoolbox}
\AtBeginEnvironment{longtable}{\footnotesize}
\AtBeginEnvironment{tabular}{\footnotesize}
% edit this one line to re-date the document; \today would drift on every rebuild
\newcommand{\propdate}{2026-09-10}
\usepackage{fancyhdr}
\newcommand{\propfoot}{%
  \lfoot{\footnotesize D7065E\,--\,Project group 5, Wassim Akachi}%
  \cfoot{\footnotesize \thepage}%
  \rfoot{\footnotesize \propdate}}
\pagestyle{fancy}\fancyhf{}
\rhead{\small D7065E\,--\,Project Proposal}
\lhead{\small Week 2 Proposal}
\propfoot
\renewcommand{\headrulewidth}{0.4pt}
\renewcommand{\footrulewidth}{0.4pt}
% \maketitle forces the "plain" style on page 1; give it the footer but no
% header, since the title block already names the document there
\fancypagestyle{plain}{%
  \fancyhf{}\propfoot
  \renewcommand{\headrulewidth}{0pt}%
  \renewcommand{\footrulewidth}{0.4pt}}
\usepackage{titlesec}
\titleformat{\section}{\Large\bfseries}{\thesection}{0.6em}{}
\titleformat{\subsection}{\large\bfseries}{\thesubsection}{0.5em}{}
\titlespacing*{\subsection}{0pt}{7pt plus 2pt minus 1pt}{3pt}
\usepackage{microtype}
\setlength{\emergencystretch}{2em}
\ifLuaTeX
  \usepackage{selnolig}  % disable illegal ligatures
\fi
\usepackage{enumitem}
\setlist[itemize]{topsep=3pt,itemsep=2pt,parsep=0pt,leftmargin=1.2em}
\setlength{\parskip}{4.2pt plus 1pt minus 1pt}

\title{Project Proposal (Week 2), Ventilation Fault Detection, Diagnosis
and Graceful Degradation}
\usepackage{etoolbox}
\makeatletter
\providecommand{\subtitle}[1]{% add subtitle to \maketitle
  \apptocmd{\@title}{\par {\large #1 \par}}{}{}
}
\makeatother
\subtitle{D7065E, Embedded Intelligence at the Edge}
\author{Wassim Akachi, Project group 5 (solo)}
\date{}

\begin{document}
\maketitle

\hypertarget{summary}{%
\subsection{Summary}\label{summary}}

\begin{longtable}[]{@{}
  >{\raggedright\arraybackslash}p{(\columnwidth - 2\tabcolsep) * \real{0.2200}}
  >{\raggedright\arraybackslash}p{(\columnwidth - 2\tabcolsep) * \real{0.7800}}@{}}
\toprule
\endhead
\textbf{Use case} & Demand-controlled ventilation that detects, diagnoses
and survives its own equipment faults \\
\textbf{Pitch} & A stuck damper and a room full of people look identical
to a CO₂ controller. A model of what the last command should have
achieved tells them apart, names the fault, and keeps the air breathable
until someone repairs it. \\
\textbf{Repository} & \url{https://github.com/eduni-pro/202608-d7065e}, in subfolder \texttt{project/} \\
\bottomrule
\end{longtable}

\hypertarget{use-case-and-why}{%
\subsection{1. Use case, and why}\label{use-case-and-why}}

The system runs demand-controlled ventilation on a BuildSim office
floor: CO₂ rises with occupancy, the controller opens a damper, the air
is diluted, the next reading falls. The project proper is the
\textbf{supervisory loop} on top, which decides whether that control loop
is still \emph{working}.

It was chosen because the fault is invisible to the controller, which
only observes CO₂: a high reading looks the same whether the damper is
stuck shut or the room is simply full. Telling those apart needs a
\emph{model} of what the last command should have achieved, so the
autonomous part is an inference problem rather than a threshold, and
fault injection and degraded operation become the use case itself rather
than an experiment bolted on at the end.

\hypertarget{scenario}{%
\subsection{2. Scenario}\label{scenario}}

A room fills up and CO₂ climbs past 1000\,ppm. The controller opens the
damper wide, but it is stuck, so the airflow never changes and CO₂ keeps
rising. The monitor saw the same command, expected the level to fall,
and instead watches the gap grow until it is too large and too
persistent to be noise. Unchanged fan power argues against a blocked
filter, a correct overnight decay argues against a drifting sensor, and
the monitor publishes a confirmed \emph{stuck damper}. The supervisor
degrades that room, over-ventilates its neighbour to hold the occupants
below the limit, and raises a work order shown on the 3D view. The
diagnosis stays latched even though the CO₂ curve now looks healthy
again.

\hypertarget{what-will-be-built}{%
\subsection{\texorpdfstring{3. What will be built \emph{(rough
sketch)}}{3. What will be built (rough sketch)}}\label{what-will-be-built}}

\begin{itemize}
\tightlist
\item
  \textbf{Physics simulator:} advances the CO₂ and occupancy models,
  applies the injected fault, and writes the resulting sensor values
  into BuildSim.
\item
  \textbf{Sensor and actuator processes:} publish CO₂ and fan power read
  from BuildSim; validate and apply damper commands back to BuildSim,
  which stays the single source of building state.
\item
  \textbf{Controller:} the inner loop, choosing a damper position per
  zone from CO₂.
\item
  \textbf{Monitor and supervisor:} the outer loop, split in two on
  purpose. The monitor predicts the response each command should
  produce, watches the residual and names the fault; it raises events
  and \textbf{never commands anything}. The supervisor acts on a
  confirmed fault, degrading the zone, compensating from a neighbour and
  issuing the work order.
\item
  \textbf{Storage and dashboard:} readings, commands, residuals and
  decisions land in a time-series store that feeds both the evaluation
  and a dashboard overlaying residual and fault state per zone on the 3D
  view.
\end{itemize}

Each piece runs as its \textbf{own process}, communicating through
defined interfaces: a message stream for readings, commands and fault
events, BuildSim's API for building state. Splitting monitor from
supervisor is deliberate, a crashing monitor must not be able to stop
ventilation, and a wrong diagnosis must not be able to command anything
by itself.

\hypertarget{simulating-the-building}{%
\subsection{4. Simulating the building}\label{simulating-the-building}}

BuildSim provides the rooms but no physics, so both the air and the
people are simulated.

\begin{itemize}
\tightlist
\item
  \textbf{Physical processes:} a well-mixed CO₂ balance per zone,
  production rising with the number of people, dilution set by the
  airflow the commanded damper position allows, and volume taken from
  the room's floor area. Commands feed back into the model, so the loop
  closes. Fan power is derived from the flow, so a blocked filter costs
  more power for less air.
\item
  \textbf{The equipment fault:} injected \emph{between} the command and
  its physical effect, a stuck damper ignores the command, a blocked
  filter scales the flow down, and a drifting sensor adds a slow offset
  at the sensor rather than in the room. This model is invented rather
  than measured, and the report will say so.
\item
  \textbf{Movement of people (occupancy):} a weekday schedule per room
  type, arrivals, a lunch dip, departures, booked meetings, empty nights
  and weekends, with randomness so no two days are alike. The empty
  nights double as the reference condition for checking the sensor: with
  nobody in the building, CO₂ \emph{must} decay to the outdoor level.
\end{itemize}

\hypertarget{the-autonomous-part}{%
\subsection{5. The autonomous part}\label{the-autonomous-part}}

There are two decisions, and only the second is interesting. The
\textbf{controller} sets each zone's damper position from CO₂ with
hysteresis and a freshness check, deliberately a plain rule.

The \textbf{monitor} is the autonomous service proper. Starting
approach: predict the CO₂ response the command should have produced,
take the residual against the measurement, and accumulate it over time
(a one-sided CUSUM) so a small persistent bias is caught long before a
raw threshold would fire. A short rule tree over residual, fan power and
overnight decay separates the three faults, with an explicit
\emph{unclassified} branch that alerts a human rather than acting. A
learned alternative will be measured against that baseline on identical
seeds; if it does not win, the report will say so and keep the
baseline.

\hypertarget{biggest-unknowns-risks}{%
\subsection{6. Biggest unknowns \& risks}\label{biggest-unknowns-risks}}

\begin{itemize}
\tightlist
\item
  The prediction wants to know occupancy, the one thing a real building
  does not. The intended answer is to judge commanded \emph{step
  responses}, whose shape depends mainly on airflow and volume rather
  than head-count, but that is unproven under realistic noise.
\item
  Whether the three fault signatures stay separable once sensor noise is
  realistic, or two of them collapse into one class.
\item
  \textbf{Masking:} once the supervisor compensates, CO₂ recovers, the
  residual shrinks, and a naive monitor clears a fault that is still
  there. Latching and occasional active probing are the intended answer;
  this is the risk most likely to bite.
\item
  The fault model and the detector are both mine, so the evaluation
  measures the detector against my own generator.
\end{itemize}

\hypertarget{first-step-rough-plan}{%
\subsection{7. First step \& rough plan}\label{first-step-rough-plan}}

The riskiest assumption is that a residual separates a stuck damper from
a busy room at all, so that is the first prototype: one zone, offline, a
step response with and without the fault, and a plot of the two
residuals. If they do not separate, the use case changes now rather than
in week 8. After that, one zone end-to-end to prove the loop closes;
monitor and supervisor split into their own processes; more zones,
storage and the dashboard; the fault-injection harness; the experiments,
detection delay, false alarms and a scaling run to the breaking point;
then report and demo.

\end{document}
